\documentclass[10pt,a4paper]{amsart}
\usepackage{amsmath,amssymb,amsthm,graphicx}
% Retained preamble for the Stacks Project chapter snapshot in this directory.
% The chapter's own preamble.tex is a copy of the shared Stacks Project preamble; the block
% below is sliced MECHANICALLY, line for line, from it. Nothing was added, removed or
% redefined; the wrapper class, the xr-hyper/hyperref wiring and the externaldocument list
% are omitted because the assembler supplies the document class.
%
% lines taken from preamble.tex: 3, 5, 9, 11, 12, 156, 157, 158, ... (51 lines total)
%
% package coverage restored from the chapter's own preamble (2026-10-13): xy, xy-2cell

\usepackage[all]{xy}
\xyoption{2cell}
\UseAllTwocells

\newenvironment{reference}{\comment}{\endcomment}
\newenvironment{slogan}{\comment}{\endcomment}
\newenvironment{history}{\comment}{\endcomment}
\theoremstyle{plain}
\newtheorem{theorem}[subsection]{Theorem}
\newtheorem{proposition}[subsection]{Proposition}
\newtheorem{lemma}[subsection]{Lemma}
\theoremstyle{definition}
\newtheorem{definition}[subsection]{Definition}
\newtheorem{example}[subsection]{Example}
\newtheorem{exercise}[subsection]{Exercise}
\newtheorem{situation}[subsection]{Situation}
\theoremstyle{remark}
\newtheorem{remark}[subsection]{Remark}
\newtheorem{remarks}[subsection]{Remarks}
\numberwithin{equation}{subsection}
\def\lim{\mathop{\mathrm{lim}}\nolimits}
\def\colim{\mathop{\mathrm{colim}}\nolimits}
\def\Spec{\mathop{\mathrm{Spec}}}
\def\Hom{\mathop{\mathrm{Hom}}\nolimits}
\def\Ext{\mathop{\mathrm{Ext}}\nolimits}
\def\SheafHom{\mathop{\mathcal{H}\!\mathit{om}}\nolimits}
\def\SheafExt{\mathop{\mathcal{E}\!\mathit{xt}}\nolimits}
\def\Sch{\mathit{Sch}}
\def\Mor{\mathop{\mathrm{Mor}}\nolimits}
\def\Ob{\mathop{\mathrm{Ob}}\nolimits}
\def\Sh{\mathop{\mathit{Sh}}\nolimits}
\def\NL{\mathop{N\!L}\nolimits}
\def\CH{\mathop{\mathrm{CH}}\nolimits}
\def\proetale{{pro\text{-}\acute{e}tale}}
\def\etale{{\acute{e}tale}}
\def\QCoh{\mathit{QCoh}}
\def\Ker{\mathop{\mathrm{Ker}}}
\def\Im{\mathop{\mathrm{Im}}}
\def\Coker{\mathop{\mathrm{Coker}}}
\def\Coim{\mathop{\mathrm{Coim}}}
\DeclareMathSymbol{\boxtimes}{\mathbin}{AMSa}{"02}
\def\QCohstack{\mathcal{QC}\!\mathit{oh}}
\def\Cohstack{\mathcal{C}\!\mathit{oh}}
\def\Spacesstack{\mathcal{S}\!\mathit{paces}}
\def\Quotfunctor{\mathrm{Quot}}
\def\Hilbfunctor{\mathrm{Hilb}}
\def\Curvesstack{\mathcal{C}\!\mathit{urves}}
\def\Polarizedstack{\mathcal{P}\!\mathit{olarized}}
\def\Complexesstack{\mathcal{C}\!\mathit{omplexes}}
\def\Pic{\mathop{\mathrm{Pic}}\nolimits}
\def\Picardstack{\mathcal{P}\!\mathit{ic}}
\def\Picardfunctor{\mathrm{Pic}}
\def\Deformationcategory{\mathcal{D}\!\mathit{ef}}
\usepackage{cleveref}
\crefname{originallemma}{original lemma}{original lemmas}
\Crefname{originallemma}{Original lemma}{Original lemmas}
\makeatletter
\newlabel{algebra-lemma-cokernel-flat}{{00O0}{source}{}{sourcepointer-1}{}}
\newlabel{algebra-lemma-cokernel-flat@cref}{{[originallemma][0][]00O0}{[1][0][]source}{}{}{}}
\makeatother
\crefname{originallemma}{original lemma}{original lemmas}
\Crefname{originallemma}{Original lemma}{Original lemmas}
\makeatletter
\newlabel{lemma-hilb-d-of-closed}{{0B97}{source}{}{sourcepointer-2}{}}
\newlabel{lemma-hilb-d-of-closed@cref}{{[originallemma][0][]0B97}{[1][0][]source}{}{}{}}
\makeatother
\crefname{originallemma}{original lemma}{original lemmas}
\Crefname{originallemma}{Original lemma}{Original lemmas}
\makeatletter
\newlabel{lemma-hilb-d-sheaf}{{0B95}{source}{}{sourcepointer-3}{}}
\newlabel{lemma-hilb-d-sheaf@cref}{{[originallemma][0][]0B95}{[1][0][]source}{}{}{}}
\makeatother
\crefname{originallemma}{original lemma}{original lemmas}
\Crefname{originallemma}{Original lemma}{Original lemmas}
\makeatletter
\newlabel{schemes-lemma-glue-functors}{{01JJ}{source}{}{sourcepointer-4}{}}
\newlabel{schemes-lemma-glue-functors@cref}{{[originallemma][0][]01JJ}{[1][0][]source}{}{}{}}
\makeatother
\title{Hilb\textasciicircum{}d of an affine scheme is representable}
\author{The Stacks Project Authors}
\date{stacks-project:pic.tex:local-snapshot-20261012; GNU FDL 1.2 or later}
\begin{document}
\maketitle
\noindent\textit{Editorial note.} Statement span is lines 410-414 of sources/stacks-pic/pinned-original-source.tex; proof span is lines 416-513, the complete author proof (96 body lines).\par\medskip
\noindent\textit{Editorial note.} Navigation stubs: 4 cited-but-not-included labels are supplied as non-load-bearing pointers with their pinned tags where the snapshot tag index has them.\par\medskip
\noindent\textit{Editorial note.} extra\_packages is empty on purpose: the retained preamble already loads xy with options and a second declaration clashes in the assembler wrapper.\par\medskip
\noindent\textit{Editorial note.} Text is retained verbatim, including the author's own typos and placeholders.\par\medskip
\noindent\textit{Editorial note.} Queue index 1 of queue\_pic.json (pb=4020, stmt 410-414, proof 416-513).\par\medskip
\noindent\textit{Editorial note.} stacks-pic wave-1 (pi01): fresh adjudication accept (2026-10-14); pool per standing queue hypercovering -> pic (SMC reassigned to Zhen per Chief 08:5x); pic.tex 27 labels/23 with proof judged 21 accept / 2 reject / 4 pb=0 env-dup rows excluded; spans cite-free and reference/slogan/history-free so no patchenv needed.\par\medskip
\noindent\textit{Allowed standard prerequisites (editorial).}\begin{itemize}
\item Explicit prerequisite boundary. The cited inputs are prerequisites; the argument in the retained span is not replaced by any cited result.
\end{itemize}\medskip
\ifdefined\hypertarget\hypertarget{sourcepointer-1}{}\fi
\noindent\textit{Source reference.} Original lemma 00O0: sources/stacks-pic/pinned-original-source.tex, source lines 513--513. Cited by the retained statement or proof but not part of the collected spans; supplied only as a navigation pointer. The argument does not depend on its content.\par\medskip
\ifdefined\hypertarget\hypertarget{sourcepointer-2}{}\fi
\noindent\textit{Source reference.} Original lemma 0B97: sources/stacks-pic/pinned-original-source.tex, source lines 341--347. Cited by the retained statement or proof but not part of the collected spans; supplied only as a navigation pointer. The argument does not depend on its content.\par\medskip
\ifdefined\hypertarget\hypertarget{sourcepointer-3}{}\fi
\noindent\textit{Source reference.} Original lemma 0B95: sources/stacks-pic/pinned-original-source.tex, source lines 282--286. Cited by the retained statement or proof but not part of the collected spans; supplied only as a navigation pointer. The argument does not depend on its content.\par\medskip
\ifdefined\hypertarget\hypertarget{sourcepointer-4}{}\fi
\noindent\textit{Source reference.} Original lemma 01JJ: sources/stacks-pic/pinned-original-source.tex, source lines 513--513. Cited by the retained statement or proof but not part of the collected spans; supplied only as a navigation pointer. The argument does not depend on its content.\par\medskip
\begin{lemma}
\label{lemma-hilb-d-An}
Let $X \to S$ be a morphism of affine schemes. Let $d \geq 0$. Then
$\Hilbfunctor^d_{X/S}$ is representable.
\end{lemma}

\begin{proof}
Say $S = \Spec(R)$. Then we can choose a closed immersion of $X$
into the spectrum of $R[x_i; i \in I]$ for some set $I$ (of sufficiently
large cardinality. Hence by Lemma \ref{lemma-hilb-d-of-closed}
we may assume that $X = \Spec(A)$ where $A = R[x_i; i \in I]$.
We will use Schemes, Lemma \ref{schemes-lemma-glue-functors} to prove the
lemma in this case.

\medskip\noindent
Condition (1) of the lemma follows from Lemma \ref{lemma-hilb-d-sheaf}.

\medskip\noindent
For every subset $W \subset A$ of cardinality $d$ we will
construct a subfunctor $F_W$ of $\Hilbfunctor^d_{X/S}$.
(It would be enough to consider the case where $W$ consists of a
collection of monomials in the $x_i$ but we do not need this.)
Namely, we will say that $Z \in \Hilbfunctor^d_{X/S}(T)$ is in $F_W(T)$
if and only if the $\mathcal{O}_T$-linear map
$$
\bigoplus\nolimits_{f \in W} \mathcal{O}_T
\longrightarrow
(Z \to T)_*\mathcal{O}_Z,\quad
(g_f) \longmapsto \sum g_f f|_Z
$$
is surjective (equivalently an isomorphism). Here for $f \in A$
and $Z \in \Hilbfunctor^d_{X/S}(T)$ we denote $f|_Z$ the pullback of $f$
by the morphism $Z \to X_T \to X$.

\medskip\noindent
Openness, i.e., condition (2)(b) of the lemma. This follows from
Algebra, Lemma \ref{algebra-lemma-cokernel-flat}.

\medskip\noindent
Covering, i.e., condition (2)(c) of the lemma. Since
$$
A \otimes_R \mathcal{O}_T =
(X_T \to T)_*\mathcal{O}_{X_T} \to (Z \to T)_*\mathcal{O}_Z
$$
is surjective and since $(Z \to T)_*\mathcal{O}_Z$ is finite locally
free of rank $d$, for every point $t \in T$ we can find a finite
subset $W \subset A$ of cardinality $d$ whose images form a basis
of the $d$-dimensional $\kappa(t)$-vector space
$((Z \to T)_*\mathcal{O}_Z)_t \otimes_{\mathcal{O}_{T, t}} \kappa(t)$.
By Nakayama's lemma there is an open neighbourhood $V \subset T$
of $t$ such that $Z_V \in F_W(V)$.

\medskip\noindent
Representable, i.e., condition (2)(a) of the lemma. Let $W \subset A$
have cardinality $d$. We claim that $F_W$ is representable by an affine
scheme over $R$. We will construct this affine scheme here, but we encourage
the reader to think it through for themselves. Choose a numbering
$f_1, \ldots, f_d$ of the elements of $W$. We will construct a universal
element $Z_{univ} = \Spec(B_{univ})$ of $F_W$ over $T_{univ} = \Spec(R_{univ})$
as follows. First, we set
$$
R_{univ} = R[c_{kl}^m, b^l, b_i^l]/\mathfrak a_{univ}
$$
for some ideal $\mathfrak a_{univ}$ which we describe below.
Here the indices $k, l, m, i$ are all ranging through $\{1, \ldots, d\}$.
Then we set
$$
B_{univ} = R_{univ}e_1 \oplus R_{univ}e_2 \oplus \ldots \oplus R_{univ}e_d
$$
as an $R_{univ}$-module with algebra structure given by the rule
$$
e_ke_l = \sum c_{kl}^m e_m
$$
The closed immersion $Z_{univ} \to X_{T_{univ}}$ is given
by the $R_{univ}$-algebra map
$$
\Psi : A \otimes_R R_{univ} \longrightarrow B_{univ}
$$
mapping $1 \otimes 1$ to $\sum b^le_l$ and $x_i$ to $\sum b_i^le_l$
(for existence see below). As for the ideal $\mathfrak a_{univ}$
we have the following prescription:
\begin{enumerate}
\item multiplication on $B_{univ}$ is commutative, i.e.,
we should have $c_{lk}^m - c_{kl}^m \in \mathfrak a_{univ}$,
\item multiplication on $B_{univ}$ is associative, i.e.,
we should have $c_{lk}^m c_{m n}^p - c_{lq}^p c_{kn}^q \in \mathfrak a_{univ}$,
\item $\sum b^le_l$ is a multiplicative $1$ in $B_{univ}$,
in other words, we should have $(\sum b^le_l)e_k = e_k$ for all $k$,
which means $\sum b^lc_{lk}^m - \delta_{km} \in \mathfrak a_{univ}$
(Kronecker delta).
\end{enumerate}
Since we've guaranteed that $B_{univ}$ is a commutative
$R_{univ}$-algebra with $1$ given by $\sum b^le_l$ the homomorphism $\Psi$
exists. The final condition is that
\begin{enumerate}
\item[(5)] $f_l$ to map to $e_l$ in $B_{univ}$. Write
$\Psi(f_l) - e_l \equiv \sum h_l^me_m$ for some
$h_l^m \in R[c_{kl}^m, b^l, b_i^l]$
then we should have $h_l^m \in \mathfrak a_{univ}$.
\end{enumerate}
With $\mathfrak a_{univ} \subset R[c_{kl}^m, b^l, b_i^l]$ generated
by the elements listed in (1) -- (5)
it is clear that $F_W$ is represented by $\Spec(R_{univ})$.
\end{proof}

\end{document}
